\begin{riskbox}[核心风险]
第一，估值拥挤。多只核心标的已在午盘成交额和涨幅上体现强情绪，若基本面兑现慢于预期，回撤会比行业需求变化更快。第二，客户集中。AI 服务器、光模块和 PCB 的订单高度依赖少数 CSP 或平台客户。第三，AIDC 运营的瓶颈在交付后：上架率、电价、折旧和运维能力比拿地和建设更重要。第四，技术路线变化可能转移利润池，例如 CPO/LPO、铜互联回潮、国产芯片架构差异或液冷方案标准变化。
\end{riskbox}

\begin{exhibitbox}[催化剂与监测阈值]
\begin{tabularx}{\textwidth}{L{3.0cm}X X}
\toprule
\textbf{监测项} & \textbf{正向阈值} & \textbf{负向阈值}\\
\midrule
服务器/交换机 & 工业富联、浪潮信息等 Q2/Q3 收入继续高增且毛利率不再下探 & 增收不增利、存货和应收显著扩张\\
光互联 & 800G/1.6T 出货、毛利率和现金流同步提升 & ASP 下行、客户砍单或库存堆积\\
AI PCB/CCL & 高多层/HDI/高速交换机板收入占比继续提升 & 扩产过快、良率或价格压力压缩毛利\\
液冷/供配电 & CDU、冷板、UPS、电力模组从认证进入批量交付 & 只有样机/认证，没有收入和利润确认\\
AIDC运营 & 新增 MW 上架率、客户粘性和经营现金流改善 & 高折旧、低利用率、融资成本上行\\
\bottomrule
\end{tabularx}
\sourcenote{analysis/risk\_framework.md；analysis/chain\_earnings\_bridge.md。}
\end{exhibitbox}

组合执行上，报告建议把“直接收入证据”放在第一权重，“估值空间”放在第二权重，“产业链位置”只放第三权重。当前价格下更适合建立跟踪清单和事件验证表，而不是无差别追高。
